\chapter{Architektur von ROS2-Systemen}

\begin{lernziele}
Nach diesem Kapitel können Sie eine einfache ROS2-Architektur für den autonomen Roboter entwerfen.
\end{lernziele}

\section{Typische Node-Struktur}
Eine sinnvolle erste Architektur besteht aus folgenden Nodes:
\begin{itemize}
\item \texttt{camera\_node}
\item \texttt{lidar\_node}
\item \texttt{slam\_node}
\item \texttt{navigation\_node}
\item \texttt{object\_detector\_node}
\item \texttt{battery\_monitor\_node}
\item \texttt{motion\_control\_node}
\end{itemize}

\section{Kommunikationsstruktur}
\begin{verbatim}
camera_node -> /camera/image_raw -> object_detector_node
lidar_node  -> /scan             -> slam_node
slam_node   -> /map              -> navigation_node
navigation_node -> /cmd_vel      -> motion_control_node
battery_monitor_node -> /battery_state -> navigation_node
\end{verbatim}

\section{Architekturentscheidungen}
Eine wichtige Entscheidung ist die Granularität. Zu viele Nodes erhöhen Kommunikationsaufwand. Zu wenige Nodes machen das System unübersichtlich. Für den Einstieg gilt: Eine klar abgegrenzte Funktion darf ein eigener Node sein.

\section{Fehlerfälle}
Eine Architektur sollte nicht nur den Normalfall zeigen. Was passiert, wenn die Kamera ausfällt? Was passiert bei leerem Akku? Was passiert, wenn keine Karte verfügbar ist?

\begin{uebungbox}
Erstellen Sie eine ROS2-Architekturskizze mit mindestens sechs Nodes und acht Topics.
\end{uebungbox}
